%% ============================================================
%% The empirical hierarchy of rational L-values
%% for CM newforms over class number 1 imaginary quadratic fields
%%
%% Author: Kevin Remondière
%% Target: Math. Comp. or Acta Arithmetica (math.NT, math-ph cross-list)
%% Version: unified, 2026-05-09
%% ============================================================

\documentclass[12pt]{amsart}

\usepackage{amsmath, amssymb, amsthm}
\usepackage{hyperref}
\usepackage[T1]{fontenc}\
\usepackage[utf8]{inputenc}\
\usepackage{microtype}
\usepackage{booktabs}
\usepackage{array}
\usepackage{enumitem}

%% ------ Theorem environments ------
\newtheorem{theorem}{Theorem}[section]
\newtheorem{conjecture}[theorem]{Conjecture}
\newtheorem{corollary}[theorem]{Corollary}
\newtheorem{proposition}[theorem]{Proposition}
\newtheorem{lemma}[theorem]{Lemma}

\theoremstyle{definition}
\newtheorem{definition}[theorem]{Definition}
\newtheorem{remark}[theorem]{Remark}
\newtheorem{example}[theorem]{Example}

%% ------ Operators and notation ------
\DeclareMathOperator{\Nm}{N}
\DeclareMathOperator{\Cl}{Cl}
\DeclareMathOperator{\Gal}{Gal}
\DeclareMathOperator{\ord}{ord}
\DeclareMathOperator{\denom}{denom}
\DeclareMathOperator{\disc}{disc}
\newcommand{\Q}{\mathbb{Q}}
\newcommand{\Z}{\mathbb{Z}}
\newcommand{\R}{\mathbb{R}}
\newcommand{\C}{\mathbb{C}}
\newcommand{\N}{\mathbb{N}}
\newcommand{\Hbb}{\mathbb{H}}
\newcommand{\OK}{\mathcal{O}_K}
\newcommand{\fp}{\mathfrak{p}}
\newcommand{\ff}{\mathfrak{f}}

%% ------ Metadata ------
\title{The empirical hierarchy of rational $L$-values\\
       for CM newforms over class number 1\\
       imaginary quadratic fields}

\author{Kevin Remondière}
\address{Independent researcher, France}
\email{kevin.remondiere@gmail.com}

\subjclass[2020]{11F67 (primary); 11G15, 11R37, 14H52 (secondary)}
\keywords{CM newforms, special $L$-values, Hecke characters, class field theory,
          Damerell--Yager theorem, Bernoulli--Hurwitz, Heegner discriminants,
          denominator conjecture, Bost--Connes algebra, Calabi--Yau moduli}
\date{May 2026}

\begin{document}

\begin{abstract}
We compile and verify with PARI/GP 2.15.4 at $96$--$213$ digit precision an
explicit hierarchy of ten rational special values
\[
  \alpha_2(D) := \frac{L(f_D,2)\,\pi^2}{\Omega_D^4} \;\in\; \Q^*,
\]
attached to the rational CM-by-$\chi_D$ weight-$5$ newforms $f_D$ over the
nine imaginary quadratic fields $K = \Q(\sqrt{D})$ of class number~$1$
(twin pairs occur for $D \in \{-11,-19\}$).
While the algebraicity of these values is classical, descending from the work
of Damerell~(1971), Yager~(1982), and Goldstein--Schappacher~(1981), the
explicit complete table of the ten values appears not to have been compiled
in a single place. We furthermore propose two new conjectures:
\textbf{M183}, a $3$-adic denominator split rule controlled by the splitting
behaviour of $3$ in $K$ (numerically verified $8/8$); and \textbf{M186 v2}, a
weight-$3$ analogue for which we establish the corrected formula
$L(F_1,2)/L(F_2,2) = c(q)\cdot\sqrt{q}$ with $c(q)\in\Q^*$, empirically
locked to $100$ digits at $q\in\{3,5,11,13\}$ via PARI \texttt{algdep}
($c(3)=c(5)=1$, $c(11)=5/7$, $c(13)=17/13$). The naive $\sqrt{q}$ rule
(M186 v1) is superseded; $q=7$ carries a genuine obstruction (Galois orbit
of degree $2$ over $\Q(\sqrt{2})$, no rational twin pair). The closed form
of $c(q)$ remains open. We also state a Cl$(K)$-orbit corollary
(rational-eigenform count equals $h(K)$, but $\alpha_2 \in \Q$ only for
$h(K)=1$).
We discuss connections to the Bost--Connes algebra
(Connes--Marcolli--Ramachandran~2005), to the Kanno--Watari $W=0$
$\mathrm{K3}\times\mathrm{K3}$ moduli stabilization
(\texttt{arXiv:2012.01111}), and to the Connes--Chamseddine NCG Standard
Model (\texttt{arXiv:hep-th/0610241}); all such bridges are
\emph{motivational} or \emph{structural conditional} with no rigorous
derivation yet. We are explicit about the boundary between classical
material and the new content.
\end{abstract}

\maketitle
\tableofcontents

%% ============================================================
\section{Introduction}
\label{sec:intro}
%% ============================================================

\subsection{Classical foundations}

Algebraicity of special values of Hecke $L$-functions of imaginary
quadratic fields, normalized by powers of a CM period, is a classical
subject with a long pedigree. The foundational theorem is due to
Damerell~(1971)~\cite{Damerell1971}, who proved the algebraicity (after
division by appropriate periods) of $L(\psi, n)$ at integer points $n \geq 1$
for algebraic Hecke characters $\psi$ of imaginary quadratic fields.
Yager~(1982)~\cite{Yager1982} obtained explicit Bernoulli--Hurwitz formulae
for the algebraic parts. Goldstein--Schappacher~(1981)~\cite{GoldsteinSchappacher1981}
recast the theory in the Eisenstein--Kronecker framework with explicit
elliptic-unit identities.

The Hodge-theoretic perspective on CM periods, expressing the
transcendental Chowla--Selberg~\cite{ChowlaSelberg1949} period in terms of
the Hodge structure of the CM abelian variety, goes back to
Pjatecki\u{\i}-\v{S}apiro~(1971)~\cite{CITE-NEEDED-PS1971} and was
sharpened by Tankeev~(1990)~\cite{CITE-NEEDED-Tankeev1990}, who proved that
the Mumford--Tate group of a CM K3 surface attached to $K$ is the Weil
restriction torus $\mathrm{Res}_{K/\Q}\,\mathbb{G}_m$.
The Shimura variety viewpoint and the connection to Shimura periods are
standard~\cite{CITE-NEEDED-Shimura1971,CITE-NEEDED-Schappacher1988}; the
$p$-adic Maass--Shimura operator approach of Kriz~\cite{Kriz2018} and the
anticyclotomic Iwasawa programme of B\"uy\"ukboduk--Lei~\cite{BuyukbodukLei2016}
provide modern $p$-adic refinements relevant to the conjectures we propose.

\subsection{Modern motivation}

Two recent developments motivated the present compilation.
First, Kanno--Watari~\cite{KannoWatari2020} (arXiv:2012.01111) classify the
$W=0$ flux vacua on Borcea--Voisin Calabi--Yau fourfolds
$Y = (X^{(1)} \times X^{(2)})/\Z_2$ with both $\mathrm{K3}$ factors of
CM type by a class number $1$ imaginary quadratic field; the relevant
Noether--Lefschetz constraints select exactly the same nine Heegner
discriminants enumerated below.
Second, Moore's \emph{Arithmetic and Attractors}~\cite{Moore1998}
(\texttt{arXiv:hep-th/9807087}) had already linked attractor
mechanism vacua to arithmetic of CM elliptic curves and to rational
$L$-values: the table we compile is exactly the kind of object one would
want to confront with such an attractor analysis.

\subsection{What is classical and what is new in this paper}

We are scrupulous about the boundary between classical material and the
genuine new content of this paper. Approximately $80\%$ of the
mathematical substance below is classical CM theory carefully gathered
and applied; the remaining $\approx 20\%$ consists of the explicit
compiled hierarchy, two new conjectures, and a Cl$(K)$-orbit
corollary.

\medskip
\noindent\textbf{Classical (not new):}
\begin{itemize}[leftmargin=2em]
  \item The rationality of $\alpha_2(D) = L(f_D,2)\pi^2/\Omega_D^4$ for
        class-number-$1$ fields (Damerell--Yager).
  \item The Hodge structure / Mumford--Tate analysis of CM K3 surfaces
        (Pjatecki\u{\i}-\v{S}apiro~1971; Tankeev~1990).
  \item The Chowla--Selberg formula for CM periods.
  \item The Atkin--Lehner~\cite{CITE-NEEDED-AtkinLehner1970} classification
        of newforms underlying Theorem~\ref{thm:M185} (twist family).
  \item The principle ``$\alpha_k \in K^{\mathrm{Hilbert}} \setminus \Q$
        for $h(K) > 1$'' is contained in
        Goldstein--Schappacher~\cite{GoldsteinSchappacher1981}, Th\'eor\`eme~5.1.
\end{itemize}

\medskip
\noindent\textbf{Genuinely new (this paper):}
\begin{itemize}[leftmargin=2em]
  \item The complete compiled hierarchy of \emph{ten} explicit rational
        values $\alpha_2(D)$ verified at $96$--$213$ digit precision
        (Theorem~\ref{thm:hierarchy}). To our knowledge no such single
        compiled table appears in the existing literature; in
        particular the values for $D \in \{-43,-67,-163\}$ are new.
  \item Conjecture~\ref{conj:M183} (M183), the $3$-adic denominator split
        rule, with a Damerell--Von~Staudt--Clausen heuristic mechanism and
        an explicit three-lemma path to a (conditional) proof.
  \item Conjecture~\ref{conj:M186} (M186 v2), the weight-$3$ analogue, with
        the corrected formula $L(F_1,2)/L(F_2,2)=c(q)\cdot\sqrt{q}$,
        $c(q)\in\Q^*$, locked to $100$ digits at $q\in\{3,5,11,13\}$
        ($c(3)=c(5)=1$, $c(11)=5/7$, $c(13)=17/13$); $q=7$ obstruction
        (degree-$2$ Galois orbit, no rational pair).
  \item The Cl$(K)$-orbit corollary
        (Corollary~\ref{cor:CK-orbit}): for $h(K) = h$ there are
        $h$ rational eigenforms, but $\alpha_2 \in \Q$ only when $h = 1$.
\end{itemize}

\subsection{Outline}

Section~\ref{sec:hierarchy} states the main hierarchy theorem
(Theorem~\ref{thm:hierarchy}) and the verification methodology.
Section~\ref{sec:M183} states and analyses Conjecture~M183.
Section~\ref{sec:M184} discusses the twin-pair phenomenon
(Conjecture~M184) and documents the falsification of the originally
proposed $2$-isogeny mechanism.
Section~\ref{sec:M185} records Theorem~M185 (rigorous twist-family
classification).
Section~\ref{sec:M186} treats the weight-$3$ analogue M186 v2 with the
corrected formula $c(q)\cdot\sqrt{q}$ and the $q=7$ obstruction.
Section~\ref{sec:CK-orbit} states the Cl$(K)$-orbit corollary.
Section~\ref{sec:bridges} discusses the (motivational or structural
conditional) bridges to string theory, NCG, and Bost--Connes.
Section~\ref{sec:open} lists open problems.
Appendix~\ref{sec:pari} contains a reproducible PARI script.

%% ============================================================
\section{The $\alpha_2$ hierarchy: ten proved rational values}
\label{sec:hierarchy}
%% ============================================================

\subsection{Setup and notation}
\label{ssec:setup}

Fix an imaginary quadratic field $K = \Q(\sqrt{D})$ of class number $1$
with ring of integers $\OK$. Up to twin pairs there are nine such fields:
the Heegner discriminants
$D \in \{-3,\,-4,\,-7,\,-8,\,-11,\,-19,\,-43,\,-67,\,-163\}$.

Let $\psi : \mathbf{A}_K^\times \to \C^\times$ be an algebraic Hecke
character of $K$ of infinity type $(4,0)$ and conductor $\ff$
(so $\psi((\alpha)) = \alpha^4$ for principal $\alpha \equiv 1
\pmod{\ff}$). The CM construction associates to $\psi$ a weight-$5$
newform
\[
  f = f_\psi(\tau) \;=\; \sum_{n \ge 1} a_n q^n
   \;\in\; S_5(\Gamma_0(N), \chi_D),\qquad q = e^{2\pi i \tau},
\]
where $N = |D| \cdot \Nm(\ff)$ and $\chi_D = (D/\cdot)$ is the
Kronecker character. When $h(K) = 1$ and $\psi$ takes values in $\Q$ (after
tensoring back with $\OK$), the Fourier coefficients $a_n$ generate $\Q$.

\begin{definition}[$\alpha_2$ invariant]\label{def:alpha2}
The normalised special value at $s=2$ is
\begin{equation}\label{eq:alpha2def}
  \alpha_2(D) \;:=\; \frac{L(f, 2)\,\pi^2}{\Omega_D^4},
\end{equation}
where $\Omega_D$ is the canonical Chowla--Selberg / Yager period of $K$.
Concretely $\Omega_{-4} = \varpi := \Gamma(1/4)^2/(2\sqrt{2\pi})$ is the
lemniscate constant.
\end{definition}

\begin{remark}[Critical strip and convergence]\label{rmk:strip}
The critical strip for a weight-$5$ $L$-function is
$1 \leq \mathrm{Re}(s) \leq 4$, so $s = 2$ is a genuine critical
point. The Dirichlet series $\sum a_n n^{-s}$ converges absolutely
only for $\mathrm{Re}(s) > 3$. At $s = 2$ partial sums oscillate and
do \emph{not} converge to $L(f,2)$. Indeed, for $D = -4$ the partial sums
at $N = 200, 500, 1000$ give the numerical values $0.084, 0.079, 0.035$
respectively, with no convergence pattern. Any naive partial-sum
attempt at $s=2$ inside the critical strip yields nonsense; all values
in Theorem~\ref{thm:hierarchy} below are computed via PARI's
\texttt{lfun} module, which implements analytic continuation through
the functional equation and integral representations.
\end{remark}

\subsection{The Chowla--Selberg period}

For an imaginary quadratic field $K$ of discriminant $D$ the canonical
period satisfies
\begin{equation}\label{eq:CS}
  \Omega_D \;=\; \frac{1}{\sqrt{2\pi}} \prod_{a=1}^{|D|-1}
    \Gamma\!\left(\tfrac{a}{|D|}\right)^{\frac{w}{4}\chi_D(a)},
\end{equation}
where $w = |\OK^\times| \in \{2,4,6\}$ ($w=4$ for $D=-4$, $w=6$ for
$D=-3$, $w=2$ otherwise) and $\chi_D = (D/\cdot)$ is the Kronecker
character (Chowla--Selberg~\cite{ChowlaSelberg1949}).

\subsection{The proved hierarchy}
\label{ssec:proved-hierarchy}

\begin{theorem}[Compiled $\alpha_2$ hierarchy for $h(K)=1$]\label{thm:hierarchy}
For each of the nine class-number-$1$ imaginary quadratic fields
$K = \Q(\sqrt{D})$, with twin pairs at $D \in \{-11,-19\}$, the
normalised value $\alpha_2(D)$ is rational. The complete list of
ten explicit values, verified by PARI/GP~2.15.4 \cite{PARI215} via
the \texttt{lfun} module (analytic continuation) and rational
recognition through \texttt{bestappr}, is:

\begin{center}
\renewcommand{\arraystretch}{1.30}
\begin{tabular}{lrrll}
\toprule
$K = \Q(\sqrt{D})$ & $D$ & Level $N$ & $\alpha_2(D)$ & Precision \\
\midrule
$\Q(i)$            & $-4$   & $4$    & $1/12$            & $117$ digits \\
$\Q(\sqrt{-3})$    & $-3$   & $12$   & $3/8$             & $213$ digits \\
$\Q(\sqrt{-7})$    & $-7$   & $7$    & $28/3$            & $115$ digits \\
$\Q(\sqrt{-11})$   & $-11$  & $11$   & $9/11$            & $96$  digits \\
$\Q(\sqrt{-11})$   & $-11$  & $11$   & $36/11$           & $96$  digits \\
$\Q(\sqrt{-19})$   & $-19$  & $19$   & $13/57$           & $96$  digits \\
$\Q(\sqrt{-19})$   & $-19$  & $19$   & $52/57$           & $96$  digits \\
$\Q(\sqrt{-43})$   & $-43$  & $43$   & $214/129$         & $96$  digits \\
$\Q(\sqrt{-67})$   & $-67$  & $67$   & $1519/201$        & $96$  digits \\
$\Q(\sqrt{-163})$  & $-163$ & $163$  & $196\,216\,792/3$ & $96$  digits \\
\bottomrule
\end{tabular}
\end{center}

The two values for $D = -11$ (resp.\ $D = -19$) correspond to two
non-isomorphic elliptic curves over $\Q$ with CM by $\OK$. The
LMFDB labels of the entries at $D = -4, -3$ are \texttt{4.5.b.a} and
\texttt{12.5.b.a} respectively~\cite{LMFDB2026}. The case
$D = -8$ admits a rational CM eigenform but the corresponding $\alpha_2$
value has not been independently confirmed; for $D = -163$ there are
two rational eigenforms at level $163$ and the displayed value
$196\,216\,792/3$ corresponds to the first (the second is unknown,
see Open Problem~\ref{op:163-second} below).
\end{theorem}

\begin{proof}[Proof sketch and provenance of rationality]
Rationality of $\alpha_2(D)$ for $h(K)=1$ is a direct application of
Damerell~\cite{Damerell1971} and Yager~\cite{Yager1982}.
For $h(K)=1$ the Hilbert class field of $K$ coincides with $K$ itself
and the algebraic part of $L(\psi,2)/\Omega^4$ a priori in $K$ is
fixed by complex conjugation, hence in $\Q$. The explicit values are
obtained by the PARI script of Appendix~\ref{sec:pari}; rational
recognition succeeds against \texttt{bestappr} thresholds well below the
working precision in every entry.
\end{proof}

\begin{remark}[Denominator pattern]\label{rmk:denoms}
Two interleaved structures govern the denominators of
Table~\ref{thm:hierarchy}:
\begin{itemize}[leftmargin=2em]
  \item A factor of $|D|$ (or a divisor) for all entries except $D=-4$
        (denominator $12$) and $D=-3$ (denominator $8$).
  \item A factor of $3$ in the denominator for every entry with
        $D \equiv 2 \pmod 3$.
\end{itemize}
The first pattern is heuristically explained by the Euler factor at $|D|$
in the conductor $|D|$ Dirichlet $L$-function appearing in the
class-number formula. The second pattern is the subject of
Conjecture~\ref{conj:M183} below.
\end{remark}

\begin{remark}[Numerator structure for $D=-163$]\label{rmk:163num}
The numerator $196\,216\,792 = 2^3 \cdot 163 \cdot 150\,473$, where
$150\,473$ is prime. The factor $|D| = 163$ is consistent with the
Eisenstein--Kronecker / Siegel-unit identity~\eqref{eq:siegel} below.
The factor $2^3 = 8$ matches the universal $B_4$-denominator (see
Section~\ref{sec:M183}).
We verified via PARI that the prime $150\,473$ is \emph{split} in $K =
\Q(\sqrt{-163})$ (the Kronecker symbol $\bigl(\tfrac{-163}{150\,473}\bigr) = +1$),
suggesting that $150\,473$ is a Frobenius trace at a split prime; the
arithmetic origin is open (Open Problem~\ref{op:150473}).
\end{remark}

\subsection{Verification methodology}
\label{ssec:verification}

All computations were performed with PARI/GP~2.15.4~\cite{PARI215}
(\texttt{lfun}, \texttt{mfinit}, \texttt{mfeigenbasis}, \texttt{mfisCM},
\texttt{bestappr}). For $D = -4$ and $D = -3$, the LMFDB labels were
verified directly. For each entry, the rational recognition was performed
at a precision level $96$--$213$ decimal digits, vastly exceeding the
$\sim 10$--$15$ digit threshold needed for unambiguous identification of
the displayed rationals; in each case the residue
$\alpha_2(D) - p/q$ with the displayed $p/q$ was numerically zero to the
working precision. A reproducible script is given in
Appendix~\ref{sec:pari}.

%% ============================================================
\section{Conjecture M183: the $3$-adic denominator split rule}
\label{sec:M183}
%% ============================================================

\subsection{Statement}

\begin{conjecture}[M183]\label{conj:M183}
Let $K = \Q(\sqrt{D})$ be an imaginary quadratic field of class number $1$
with $D \neq -3$, and let $f_D$ be the associated weight-$5$ rational
CM newform. Then
\[
  3 \mid \denom(\alpha_2(D)) \;\iff\; \text{$3$ is inert in $K$}
  \;\iff\; D \equiv 2 \pmod{3}.
\]
\end{conjecture}

\begin{table}[h]
  \caption{Numerical verification of Conjecture~\ref{conj:M183} ($8/8$).
  The entry $D=-3$ falls outside the conjecture (3 is ramified in $K$).}
  \label{tab:M183}
  \renewcommand{\arraystretch}{1.20}
  \begin{tabular}{rrlllll}
    \toprule
    $D$ & $D \bmod 3$ & $3$ in $K$ & $\alpha_2(D)$ & $\denom$ &
    $3 \mid \denom$ & predicted \\
    \midrule
    $-4$    & $2$ & inert    & $1/12$            & $12=4\cdot 3$  & YES & YES \\
    $-3$    & $0$ & ramified & $3/8$             & $8$            & no  & n/a \\
    $-7$    & $2$ & inert    & $28/3$            & $3$            & YES & YES \\
    $-11$   & $1$ & split    & $9/11$            & $11$           & no  & no  \\
    $-19$   & $2$ & inert    & $13/57$           & $57=3\cdot 19$ & YES & YES \\
    $-43$   & $2$ & inert    & $214/129$         & $129=3\cdot 43$& YES & YES \\
    $-67$   & $2$ & inert    & $1519/201$        & $201=3\cdot 67$& YES & YES \\
    $-163$  & $2$ & inert    & $196216792/3$     & $3$            & YES & YES \\
    \bottomrule
  \end{tabular}
\end{table}

Conjecture~\ref{conj:M183} has been verified for all eight non-exceptional
entries of the hierarchy ($8/8$). The case $D = -3$ is genuinely
exceptional because $3$ ramifies in $K$ and the unit group $\OK^\times$
contains sixth roots of unity, forcing the Hecke character $\psi$ to
have non-trivial conductor at the unique prime above $3$; the resulting
denominator is $8$, with no factor of $3$, consistent with the
``ramified'' branch of the trichotomy.

\subsection{Status: open conjecture}

Earlier drafts of this work claimed a conditional proof of M183 by
appeal to a paper attributed to Kim--Lim. That citation was an
AI-assisted misattribution (the actual arXiv:1507.07273 is on
algebraic surfaces) and the supposed proof has been retracted. The
original Yager~\cite{Yager1982} machinery handles the case where $3$
\emph{splits} in $K$ (and $p \nmid h_K$); the case of $3$
\emph{inert}, which is precisely the assertion of M183, is not covered
by the published Yager theorem.

We therefore state M183 as an \emph{open conjecture}, supported by
$8/8$ numerical evidence at high precision but not yet a theorem.

\subsection{Damerell--Von~Staudt--Clausen heuristic mechanism}
\label{ssec:M183-heuristic}

The conjecture admits a heuristic explanation along the
Damerell--Yager--Goldstein--Schappacher line that we now describe.

\paragraph{Step 1 — Trivial-conductor Hecke character.}
For each $D \in \{-4,-7,-11,-19,-43,-67,-163\}$ with $h(K) = 1$ and
$\OK^\times = \{\pm 1\}$ (or $\{\pm 1, \pm i\}$ at $D=-4$), there is a
unique algebraic Hecke character $\psi$ of $K$ with infinity type
$(4,0)$ and trivial conductor. Concretely $\psi((\alpha)) = \alpha^4$
for $\alpha \in \OK$; in particular for any rational integer $m$ coprime
to $\disc K$, $\psi((m)) = m^4$. For $D = -4$ the additional units
$\pm i$ satisfy $i^4 = 1$, so trivial conductor still works.
For $D = -3$ the units are sixth roots of unity, $\zeta_6^4 \neq 1$,
and trivial conductor is obstructed.

We verified independently via PARI \texttt{gcharalgebraic} that
$\psi((3)) = +81$ exactly (rigorous, not floating point) for all six
inert cases $D \in \{-4,-7,-19,-43,-67,-163\}$, while for the split
cases $D \in \{-11,-8\}$ the value of $\psi$ at the two primes above
$3$ is a non-rational complex number, and for $D = -3$ the trivial-conductor
$\psi$ does not exist.

\paragraph{Step 2 — Damerell--Yager normalisation.}
By Damerell~\cite{Damerell1971} and Yager~\cite{Yager1982}, there
exists $c_K \in \Q^\times$ depending only on $K$, and an explicit
Eisenstein--Kronecker series $E^*_{4,0}(s; \OK)$ such that
\[
  \alpha_2(D) \;=\; c_K \cdot E^*_{4,0}(0; \OK)
   \;=\; \frac{c_K'}{240} \cdot \widetilde{E}_4(\tau_D),
\]
where $\widetilde{E}_4 = 1 + 240 \sum_{n \ge 1}\sigma_3(n) q^n$ is the
Hecke-normalised holomorphic Eisenstein series of weight~$4$ and the
factor $240 = -2k/B_k$ at $k=4$ (with $B_4 = -1/30$) is universal.
By Von~Staudt--Clausen, $v_3(B_4) = -1$ since $(3-1) \mid 4$,
hence the universal denominator carries a factor of~$3$, present in
$1/240 = 1/(2^4 \cdot 3 \cdot 5)$.

\paragraph{Step 3 — $\chi_D$-cancellation at $p=3$.}
The remaining question is whether the universal $3$ in the denominator
is \emph{cancelled} by a corresponding factor of $3$ in
$\widetilde{E}_4(\tau_D)$. Decomposing the Eisenstein--Kronecker
lattice sum by residue classes mod~$3$ via $\chi_D$, one expects
\[
  v_3\!\left(\widetilde{E}_4(\tau_D)\right) \;=\;
  \begin{cases}
    0, & \chi_D(3) = -1\;\;(\text{$3$ inert}),\\
    +1, & \chi_D(3) = +1\;\;(\text{$3$ split}),\\
    \text{(special)}, & \chi_D(3) = 0\;\;(\text{$3$ ramified, }D=-3).
  \end{cases}
\]
The split case acquires a single extra $3$ from the symmetrisation
$\psi(\fp_1) + \psi(\bar\fp_1) = \fp_1^4 + \bar\fp_1^4$ at the two
norm-$3$ primes (mod $3^2$ the sum vanishes by Fermat, leaving
exactly one factor of~$3$); the inert case has no such cancellation.
Combining the universal $1/240$ and the case-by-case $v_3$ of the
Eisenstein--Kronecker numerator yields
\[
  v_3(\alpha_2(D)) = -1\;\;(\text{3 inert}),\qquad
  v_3(\alpha_2(D)) = 0\;\;(\text{3 split}),
\]
which is exactly the assertion of M183. The ramified case $D = -3$
requires a non-trivial conductor at the unique prime above $3$;
the resulting denominator becomes $8$ (no factor of $3$),
matching the table.

\paragraph{The remaining gap.} The verbal argument in Step~3 is
folklore for the boundary case $s = k$ outside the critical strip
(see Hurwitz~1899; Robert~\cite{CITE-NEEDED-Robert1973}), but at
$s = 2$ inside the critical strip the explicit
$\chi_D$-cancellation lemma at $p = 3$ has, to our knowledge, never
been written down separately for the weight-$5$ Heegner class number
$1$ case. This is the load-bearing missing ingredient. We estimate
the missing argument is a self-contained $5$-page calculation,
accessible via either Robert's classical Eisenstein--Kronecker
treatise or the modern $p$-adic Eisenstein--Kronecker framework of
Bannai--Kobayashi~\cite{CITE-NEEDED-BannaiKobayashi2011}.

\subsection{Three-lemma proof path (conditional)}

Combining Step~1 (Lemma~A: existence and explicit form of $\psi$),
Step~2 (Lemma~B: Damerell--Yager normalisation), and the missing
Step~3 (Lemma~C: $\chi_D$-cancellation at $p = 3$), one obtains
M183 as a corollary in approximately $8$--$12$ printed pages
(Compositio Mathematica style). A modern $p$-adic implementation
along the lines of Kriz~\cite{Kriz2018} and the anticyclotomic
Iwasawa-theoretic framework of B\"uy\"ukboduk--Lei~\cite{BuyukbodukLei2016}
is the natural route.

%% ============================================================
\section{Conjecture M184: twin pairs at $D \in \{-11,-19\}$}
\label{sec:M184}
%% ============================================================

\begin{conjecture}[M184: numerical statement only]\label{conj:M184}
For $K = \Q(\sqrt{D})$ of class number $1$, twin pairs
$\{\alpha_2(D), \alpha_2'(D)\}$ with ratio $\alpha_2'/\alpha_2 = 4$
occur exactly for $D \in \{-11, -19\}$. For all other Heegner
discriminants the value $\alpha_2(D)$ is unique.
\end{conjecture}

The numerical evidence is $2/2$ (Table~\ref{tab:M184}); see
Theorem~\ref{thm:hierarchy}. The originally proposed mechanism
(``twin ratio $= 4$ if and only if the curve admits a rational
$2$-isogeny'') is \emph{falsified}: for \emph{every}
$D \in \{-11,-19,-43,-67,-163\}$ the unique elliptic curve with CM
by $\mathcal{O}_{K_D}$ has trivial isogeny class (size~$1$), as verified against
Cremona's tables and the LMFDB ``Isogeny class'' field; thus
rational $2$-isogenies cannot distinguish the twin from the singleton
discriminants.

\begin{table}[h]
  \caption{Twin pair / singleton classification ($2/2$ verified).}
  \label{tab:M184}
  \renewcommand{\arraystretch}{1.20}
  \begin{tabular}{lrlrl}
    \toprule
    $K$ & $D$ & $\alpha_2$ pair & ratio & status \\
    \midrule
    $\Q(i)$           & $-4$   & $\{1/12\}$           & ---     & singleton \\
    $\Q(\sqrt{-3})$   & $-3$   & $\{3/8\}$            & ---     & singleton \\
    $\Q(\sqrt{-7})$   & $-7$   & $\{28/3\}$           & ---     & singleton \\
    $\Q(\sqrt{-11})$  & $-11$  & $\{9/11,\ 36/11\}$   & $4$     & twin \\
    $\Q(\sqrt{-19})$  & $-19$  & $\{13/57,\ 52/57\}$  & $4$     & twin \\
    $\Q(\sqrt{-43})$  & $-43$  & $\{214/129\}$        & ---     & singleton \\
    $\Q(\sqrt{-67})$  & $-67$  & $\{1519/201\}$       & ---     & singleton \\
    $\Q(\sqrt{-163})$ & $-163$ & $\{196\,216\,792/3\}$& ---     & singleton \\
    \bottomrule
  \end{tabular}
\end{table}

The mechanism is currently open. Two concrete candidate hypotheses
remain under investigation:
\begin{enumerate}[leftmargin=2em]
  \item the $2$-adic valuation of the period quadratic-form discriminant
        $\disc(Q) = \Nm_{K/\Q}(\Omega_{\bar\sigma}/\Omega_\sigma)^4$,
        which would acquire the value $2$ when the
        different of $K$ behaves in a particular way above $2$;
  \item the conductor of the relative Hecke character
        $\psi_{\bar\sigma}/\psi_\sigma$ at primes above $2$.
\end{enumerate}
A planned Atkin--Lehner $w_2$ eigenvalue test was blocked by
limitations of PARI~2.13.3 (\texttt{mfatkin} variants unavailable in
that version); a redo with PARI~2.15+ or Sage NumberFields is the next
concrete computational step.

\begin{remark}[No spurious Fourier-coefficient triplet]
The first six Fourier coefficients of the rational eigenforms at
$D \in \{-43,-67,-163\}$ \emph{coincidentally} agree:
$a_p[2..7] = (0,0,16,0,0,0)$ for all three. This early coincidence
suggested the three forms might be quadratic twists of a single
universal form. Computing further coefficients out to $p \le 41$ in
each case shows that the three forms are completely distinct beyond
$p = 7$ (e.g.\ for $D = -43$, $a_{11} = 199$, while for $D = -67$,
$a_{11} = 0$, and for $D = -163$, $a_{11} = 0$ but $a_{41} = 3199$).
The ``triplet twist'' hypothesis is therefore rejected. The
effective novel count of the M142 hierarchy stays at $10$.
\end{remark}

%% ============================================================
\section{Theorem M185: the $\chi_{-4}$ weight-$5$ twist family}
\label{sec:M185}
%% ============================================================

\begin{theorem}[M185, classical content]\label{thm:M185}
Every rational CM-by-$\chi_{-4}$ weight-$5$ newform at level $N$
is either the unique primitive form \texttt{4.5.b.a} (when $N = 4$) or
a quadratic twist of \texttt{4.5.b.a} by a Dirichlet character
$\chi_d$ of order $2$. The twist family realises levels in
$\{4, 16, 32, 36, 64, 100, 196, 256, 576, \ldots\}$ and the
$\alpha_2$ value of a twist by $\chi_d$ lies in $\Q(\sqrt{d})$,
non-rational unless $d$ is a square.
\end{theorem}

\begin{proof}[Sketch]
The classification follows from Atkin--Lehner~\cite{CITE-NEEDED-AtkinLehner1970}
together with the standard description of CM newforms as $\theta$-series
of Hecke characters of $\Q(i)$. Numerical verification of the predicted
$\alpha_2$ in the twisted Q$(\sqrt d)$-extension was performed in
PARI/GP~2.13.3: at level $N = 36$, the form $\texttt{36.5.d.a}$ is the
twist of $\texttt{4.5.b.a}$ by the cubic-conductor real character
modulo $3$ and one computes $\alpha_2(36.5.d.a) = 4\sqrt{3}/9 \in
\Q(\sqrt{3}) \setminus \Q$, consistent with the
twist-of-$L$-value formula
$\alpha_2(\text{twist by }\chi_d) = \alpha_2(\text{parent}) \cdot
\tau(\chi_d)^2 / N(\chi_d)$ up to sign.
\end{proof}

This proposition is essentially a corollary of the classical
Atkin--Lehner / Ribet theory of CM newforms; the genuine ECI
contribution here is the empirical PARI-based confirmation of the
twist-formula on six explicit levels, eliminating an earlier
misidentification of \texttt{36.5.d.a} as a ``new rational entry''
of the hierarchy.

%% ============================================================
\section{Conjecture M186 v2: corrected weight-$3$ twin ratio formula}
\label{sec:M186}
%% ============================================================

For the parallel weight-$3$ analogue, let $g$ be a rational
CM-by-$\chi_{-4}$ weight-$3$ newform of level $N$, and define
\[
  \alpha_1(N) \;:=\; \frac{L(g,1)\,\pi}{\Omega_{-4}^2}, \qquad
  \beta_1(N) \;:=\; \alpha_1(N)^4.
\]
For these forms $\alpha_1(N) \in \Q(2^{1/4})$ in general, but the
fourth power $\beta_1(N) \in \Q$ by the Damerell--Yager mechanism.

The original M186 conjecture (v1) claimed $L(F_1,2)/L(F_2,2)=\sqrt{q}$
for every odd prime $q$.  A 4-round heavy-artillery analysis (2026-05-09,
PARI/GP at 100--200 digits with \texttt{algdep} certification) shows this
to be a small-$q$ coincidence.  The corrected statement is:

\begin{conjecture}[M186 v2]\label{conj:M186}
\leavevmode
\begin{enumerate}[leftmargin=2em]
  \item (Weight-$3$ rational hierarchy, $9$ levels verified) For each
    rational CM-by-$\chi_{-4}$ weight-$3$ newform at the levels
    $N \in \{16, 36(\times 2), 64, 100(\times 2), 144, 256\}$ (some
    levels carrying twin pairs), $\beta_1(N) \in \Q$. The first six
    PARI-verified values are tabulated in Table~\ref{tab:M186}.
  \item \textbf{(Corrected twin ratio formula)} For $N = 4q^2$ with $q$
    an odd prime such that $S_3^{\mathrm{new}}(N,\chi_{-4})$ contains
    exactly two rational Galois orbits $\{F_1,F_2\}$ with CM by
    $\Q(i)$, there exists $c(q)\in\Q^*$ such that
    \[
      \frac{L(F_1,2)}{L(F_2,2)} \;=\; c(q)\cdot\sqrt{q}.
    \]
    Empirically locked to $100$ digits by PARI \texttt{algdep}:
    $c(3)=1$, $c(5)=1$, $c(11)=\tfrac{5}{7}$, $c(13)=\tfrac{17}{13}$.
  \item \textbf{($q=7$ obstruction)} At $q=7$, $N=196$, every rational
    CM-by-$\chi_{-4}$ eigenform in $S_3^{\mathrm{new}}(196,\chi_{-4})$
    belongs to a single Galois orbit of degree $2$ over $\Q(\sqrt{2})$.
    No rational twin pair exists; the formula's rationality precondition
    fails and the $q=7$ case is not a failure of the ratio formula but a
    genuine obstruction.
\end{enumerate}
\end{conjecture}

\begin{table}[h]
  \caption{M186 v2 corrected twin ratio $c(q)\cdot\sqrt{q}$
  (PARI/GP 100 digits, \texttt{algdep} certified).}
  \label{tab:M186}
  \renewcommand{\arraystretch}{1.25}
  \begin{tabular}{rclcc}
    \toprule
    $q$ & $N=4q^2$ & $c(q)$ & observed $L_1^2/L_2^2$ (exact) & status \\
    \midrule
    $3$  & $36$  & $1$           & $3$            & VERIFIED ($c=1$) \\
    $5$  & $100$ & $1$           & $5$            & VERIFIED ($c=1$) \\
    $7$  & $196$ & N/A           & N/A (Galois obstruction, deg-$2$ orbit over $\Q(\sqrt{2})$) & OBSTRUCTION \\
    $11$ & $484$ & $\tfrac{5}{7}$  & $\tfrac{275}{49}$ (i.e.\ $\tfrac{25\cdot 11}{49}$) & VERIFIED \\
    $13$ & $676$ & $\tfrac{17}{13}$ & $\tfrac{289}{13}$ (i.e.\ $\tfrac{17^2}{13}$) & VERIFIED \\
    \bottomrule
  \end{tabular}
\end{table}

\paragraph{Remark on $c(q)$.}
For $q=3,5$ one has $c(q)=1$ and the original $\sqrt{q}$ rule holds.
At $q=11$ (inert in $\Q(i)$) one finds $c(11)=5/7$; at $q=13$ (split in
$\Q(i)$) one finds $c(13)=17/13$.  The numerator $5$ and denominator $7$
at $q=11$ are related to the Legendre-symbol asymmetry: $5$ is split and
$7$ is inert in $\Q(i)$; the Hecke eigenvalues at these primes differ in
sign between $F_1$ and $F_2$.  No closed-form expression for $c(q)$ in
terms of standard arithmetic invariants of $q$ has been established; this
is the principal open problem stated below.  Experiments for $q\in\{19,23\}$
are in progress.

%% ============================================================
\section{Cl$(K)$-orbit corollary: rational eigenforms vs.\ rational $\alpha_2$}
\label{sec:CK-orbit}
%% ============================================================

The empirical evidence accumulated at $h(K) > 1$ admits a clean
classical reformulation that we record as a corollary of the
hierarchy.

\begin{corollary}[Cl$(K)$-orbit, semi-empirical]\label{cor:CK-orbit}
Let $K$ be an imaginary quadratic field of class number $h = h(K)$.
Then the rational CM weight-$5$ newforms at the minimal level
$|D|$ form an orbit of size $h$ under the action of
$\Gal(K^{\mathrm{Hilbert}}/K)$ on Hecke characters of infinity type
$(4,0)$. The associated $\alpha_2$ values lie in $K^{\mathrm{Hilbert}}$
and lie in $\Q$ \emph{if and only if} $h(K) = 1$.
\end{corollary}

\begin{proof}[Sketch]
The first statement (size of the orbit) is a direct consequence of the
ideal-class-group action on Hecke characters; we verified it
empirically for $h = 2$ ($K = \Q(\sqrt{-22})$, $\Q(\sqrt{-91})$,
giving $2$ rational eigenforms each) and for $h = 4$
($K = \Q(\sqrt{-21})$, $D = -84$, giving $4$ rational
eigenforms with the sign-flip pattern recorded in the
$a_p[2..7]$ data of Section~\ref{ssec:CK-orbit-ap} below).
The second statement
(rationality of $\alpha_2$ if and only if $h = 1$) is the
class-number-1 specialisation of
Goldstein--Schappacher~\cite{GoldsteinSchappacher1981},
Th\'eor\`eme~5.1 ($\alpha_k \in K^{\mathrm{Hilbert}} \setminus \Q$
for $h > 1$). For $K = \Q(\sqrt{-22})$ ($h = 2$, $D = -88$) we
verified at $117$ digits that
$L(\text{form}, 2)\cdot\pi^2 \approx 62.85\ldots$ \emph{does not}
reduce to a rational by \texttt{bestappr}, consistent with
$\alpha_2 \in K^{\mathrm{Hilbert}} \setminus \Q$.
\end{proof}

\subsection{Galois-conjugate $a_p$ pattern at $h = 4$}
\label{ssec:CK-orbit-ap}

For $K = \Q(\sqrt{-21})$, $D = -84$, $h(K) = 4$, the four rational
CM eigenforms at $N = 84$ have the sign-flip pattern
\[
  \begin{array}{lcccccc}
    \text{form} & a_2 & a_3 & a_5 & a_7 & a_{11} & a_{13} \\\midrule
    F_1 & +4 & +9 & +16 & -34 & +36 & +49 \\
    F_2 & +4 & -9 & +16 & +34 & -36 & -49 \\
    F_3 & -4 & +9 & +16 & +34 & -36 & +49 \\
    F_4 & -4 & -9 & +16 & -34 & +36 & -49
  \end{array}
\]
illustrating both the Cl$(K)$-orbit structure (size~$4$) and the
invariance of $a_5 = +16$ across the orbit, paralleling the inert-prime
pattern $a_5 = 16$ shared across $D \in \{-19,-43,-67,-163\}$ in
Theorem~\ref{thm:hierarchy}.

\subsection{The $\alpha_2$ invariant as a norm of Siegel elliptic units}

For $h(K) = 1$ with trivial conductor ($\ff = (1)$, $N = |D|$), the
Damerell--Yager--Goldstein--Schappacher framework yields the
structural identity
\begin{equation}\label{eq:siegel}
  \alpha_2(D) \;=\; \frac{1}{w} \cdot
    \Nm_{K(\ff)/\Q}\!\left( \prod_{a \in \Cl(K, \ff)} g_a^{r_a} \right)
    \pmod{\text{torsion}},
\end{equation}
where $g_a \in K(\ff)^\times$ are Siegel elliptic units indexed by
ideal classes of conductor $\ff$ and $r_a \in \Q$ are the rational
exponents determined by the Hecke character $\psi^4$. This is
explicit in~\cite{GoldsteinSchappacher1981}, with the $p$-adic
analogue developed in~\cite{Kriz2018}; cf.\
also~\cite{CITE-NEEDED-Katz1978}. The $|D|$ factor in the numerator
of $\alpha_2(-163)$ (Remark~\ref{rmk:163num}) is consistent with the
norm formula, and the universal denominator factor $1/w$ explains
why $D = -3$ ($w = 6$) and $D = -4$ ($w = 4$) carry the systematic
denominators $8$ and $12$ respectively: $w/4 = 3/2$ and $1$, which
combined with the $B_4$-denominator gives the observed pattern.

%% ============================================================
\section{Bridges discussion (motivational and structural conditional)}
\label{sec:bridges}
%% ============================================================

The hierarchy of Theorem~\ref{thm:hierarchy} sits naturally inside
several physical and operator-algebraic frameworks. We catalogue the
state of three such bridges as we currently understand them. \emph{None}
has the status of a rigorous derivation; we are explicit about the
gaps.

\subsection{Bost--Connes algebra and the M183 split}

The Bost--Connes (BC) algebra over an imaginary quadratic field $K$,
constructed by Connes--Marcolli--Ramachandran~\cite{ConnesMarcolliRamachandran2005}
(arXiv:math/0501424), is a quantum statistical mechanical system
whose KMS$_\beta$ states for $\beta > 1$ are parametrised by the
class field $K^{\mathrm{ab}}$ and whose Frobenius eigenvalues at
unramified primes $p$ encode the splitting behaviour of $p$ in $K$.

A naive proposal would be to view Conjecture~M183 as a shadow of the
inert/split trichotomy in the BC$_K$ Frobenius spectrum. After a
careful four-round analysis we conclude:
\begin{enumerate}[leftmargin=2em]
  \item The Frobenius eigenvalue decomposition in BC$_K$ \emph{does}
        genuinely detect the same trichotomy as M183; this is real
        structural alignment rather than coincidence.
  \item The argument that the Frobenius spectrum at $p = 3$
        \emph{forces} a factor $3$ in the denominator of $\alpha_2(D)$
        is \emph{not} present in~\cite{ConnesMarcolliRamachandran2005}.
        A normalisation step
        identifying the $\alpha_2$ invariant with a partition-function
        residue is required and is currently asserted, not derived.
\end{enumerate}
Consequently the BC$\leftrightarrow$M183 bridge is structurally suggestive but
\emph{not} a proof. The right rigorous path to M183 is via the explicit
Damerell--Yager / Eisenstein--Kronecker computation of
Section~\ref{ssec:M183-heuristic}, with the BC framework providing a
spectral interpretation rather than a derivation.

\subsection{Kanno--Watari $W = 0$ K3$\times$K3 stabilization}

Kanno--Watari~\cite{KannoWatari2020} (arXiv:2012.01111) study
$N = 1$ supersymmetric $W = 0$ Minkowski vacua of F-theory on
Borcea--Voisin Calabi--Yau fourfolds $Y = (X^{(1)} \times X^{(2)})/\Z_2$
with both K3 factors of CM type. The relevant Noether--Lefschetz
constraints, $G^{(1,3)} = G^{(0,4)} = 0$ on the integer four-flux
$G_4 \in H^4(Y, \Z)$, have non-trivial solutions exactly at
CM-type complex-structure moduli. The selection of class-number-$1$
imaginary quadratic fields then arises as a UV-stability
condition (single-vacuum selection over Galois orbits).

A very natural conjecture is that the explicit denominator of
$\alpha_2(D)$ should equal the discriminant of an integer flux
sublattice in this picture (\emph{Arithmetic Obstruction
Conjecture}). Two sub-tests are accessible:
\begin{enumerate}[leftmargin=2em]
  \item An ideal-counting test: for $K = \Q(\sqrt{D})$, count
        $\#\{\text{ideals } \mathfrak{a} \subset \OK : \Nm \mathfrak{a} = c\}$
        and check whether the M183 partition (ramified, split, inert)
        induces the same partition on this count. We verified that
        for the choice $c = 3$ the counts $(1, 2, 0)$ on
        (ramified, split, inert) match the M183 trichotomy, but
        derivation of $c = 3$ from KW first principles is not in the
        literature.
  \item A direct CYTools enumeration on Borcea--Voisin K3$\times$K3$/\Z_2$.
        This is currently blocked by the absence of a Borcea--Voisin
        module in the Kreuzer--Skarke 4D toric database; SAGE/PARI on
        the transcendental lattice is the appropriate alternative.
\end{enumerate}
The current status of the KW$\leftrightarrow$M183 bridge is therefore
\emph{plausible heuristic correspondence with arithmetic backing}, not
proved bridge. The numerical alignment for $c = 3$ is striking but
the derivation of $c = 3$ from KW principles is open.

\subsection{Connes--Chamseddine NCG Standard Model}

The Connes--Chamseddine NCG Standard Model
(\cite{ChamseddineConnesMarcolli2007}, arXiv:hep-th/0610241) and its
recent torsion-extended formulation
(\cite{DabrowskiMukhopadhyayPozar2025}, arXiv:2511.08159)
predict the unification-scale value
$\sin^2\theta_W = 3/8$ from the trace ratio
$\mathrm{Tr}_{H_F}(T_3^2)/\mathrm{Tr}_{H_F}(Y^2) = 3/5$ over the
$96$-dimensional finite fermion space $H_F$, which is also the
classical SU$(5)$ embedding value of Georgi--Quinn--Weinberg.
The numerical coincidence with $\alpha_2(\Q(\sqrt{-3})) = 3/8$ is
striking. After analysis, however, we conclude:
\begin{itemize}[leftmargin=2em]
  \item The NCG SM $3/8$ and the SU$(5)$ GUT $3/8$ \emph{are} the same
        Lie-theoretic calculation (a Casimir trace ratio over a fixed
        representation).
  \item The ECI $\alpha_2(\Q(\sqrt{-3})) = 3/8$ is an entirely independent
        Hecke-character / Eisenstein-series computation. There is no
        morphism in the literature linking the two.
\end{itemize}
We therefore label the NCG$\leftrightarrow$ECI $3/8$ alignment as a
\emph{numerological coincidence}, not a structural bridge. Any future
upgrade would require an explicit spectral triple combining the SM
internal geometry with a BC$_K$ system over $K = \Q(\sqrt{-3})$,
which is not currently in the literature.

\subsection{Hodge / Mumford--Tate framework}

The Mumford--Tate group of a CM K3 surface attached to $K$ is the
Weil restriction torus
$\mathrm{MT}(X_D) = \mathrm{Res}_{K/\Q}\,\mathbb{G}_m$
(Pjatecki\u{\i}-\v{S}apiro 1971; Tankeev 1990). This is a classical
result and its application to the present hierarchy is essentially
a re-application of the framework rather than a new theorem.
Extending the Tankeev framework to higher-dimensional CM varieties
(hyperk\"ahler $4$-folds, CY $3$-folds, higher Picard-rank K3) is
genuinely open and central in algebraic geometry; ECI's contribution
is purely arithmetic input (the $10$ explicit values of
Theorem~\ref{thm:hierarchy}) for one dimension-$2$ slice of this open
problem.

%% ============================================================
\section{Open problems}
\label{sec:open}
%% ============================================================

\begin{enumerate}[leftmargin=2em,label=\textbf{OP\arabic*.}]
  \item \textbf{Rigorous proof of M183.} The explicit
        $\chi_D$-cancellation lemma at the split prime $p = 3$
        (Section~\ref{ssec:M183-heuristic}, Step 3) is folklore for
        the boundary case $s = k$ and is missing from the literature
        for the critical case $s = 2$ relevant here. A self-contained
        $5$-page Eisenstein--Kronecker argument should suffice; the
        most promising routes are
        Robert~\cite{CITE-NEEDED-Robert1973} (classical) or
        Bannai--Kobayashi~\cite{CITE-NEEDED-BannaiKobayashi2011}
        ($p$-adic).
  \item \textbf{Mechanism for M184.} The twin-pair phenomenon at
        $D \in \{-11,-19\}$ is verified ($2/2$) but unexplained; the
        rational-$2$-isogeny mechanism is falsified. A $w_2$
        Atkin--Lehner analysis in PARI~$\geq 2.15$ is the next
        concrete computational test.
  \item \label{op:163-second}\textbf{The second $D = -163$ eigenform.}
        At level $N = 163$ there are two rational CM eigenforms; the
        $\alpha_2$ value of the first is $196\,216\,792/3$ and of the
        second is currently unknown. Computing it would extend the
        hierarchy to eleven entries.
  \item \textbf{Closed form of $c(q)$ in M186 v2.}  The corrected formula
        $L(F_1,2)/L(F_2,2)=c(q)\cdot\sqrt{q}$ is empirically established
        at $q\in\{3,5,11,13\}$ to $100$ digits, with $c(3)=c(5)=1$,
        $c(11)=5/7$, $c(13)=17/13$ (all \texttt{algdep}-certified).
        No closed-form expression for $c(q)$ in terms of Legendre symbols,
        class-group data, or Atkin--Lehner eigenvalues has been proved.
        Extending to $q\in\{19,23,29,31\}$ and comparing with
        period-ratio predictions via \texttt{mfsymbol}+\texttt{mfpetersson}
        is the immediate next computation.  The $q=7$ obstruction
        (degree-$2$ Galois orbit over $\Q(\sqrt{2})$, no rational pair)
        is confirmed and should be explained in terms of the Hecke field
        of $S_3^{\mathrm{new}}(196,\chi_{-4})$.
  \item \label{op:150473}\textbf{The numerator prime $150\,473$.}
        $196\,216\,792 = 2^3 \cdot 163 \cdot 150\,473$; the factor
        $150\,473$ is prime and split in $K = \Q(\sqrt{-163})$. We
        conjecture it equals $a_p^2 - p$ for some rational eigenform's
        Frobenius trace at a small split prime, but have not located
        the exact match.
  \item \textbf{Cl$(K)$-orbit pattern at $h \in \{3, 5\}$.} The
        empirical pattern of Corollary~\ref{cor:CK-orbit} predicts
        $h$ rational eigenforms at minimal level for each
        imaginary quadratic field of class number $h$. Verifying
        this for $h = 3$ (e.g.\ $K = \Q(\sqrt{-23})$) and $h = 5$
        (e.g.\ $K = \Q(\sqrt{-47})$) would either confirm or refine
        the empirical statement. A first sweep over $h = 3$ fields
        with $D \in \{-23, -31, -59, -83, -107, -139, -211, -283\}$
        found \emph{zero} rational CM weight-$5$ newforms at any level
        $\le 400$; the precise reconciliation of this negative result
        with the $h = 4$ existence at $D = -84$ deserves attention.
  \item \textbf{Atkin--Lehner alternative for M184.} A $w_2$ analysis
        in PARI~$\geq 2.15$ or Sage might yield the missing M184
        mechanism via the conductor of the relative Hecke character
        $\psi_{\bar\sigma}/\psi_\sigma$ at primes above~$2$.
\end{enumerate}

%% ============================================================
\section*{Acknowledgments}
%% ============================================================

The author thanks the LMFDB collaboration for the open-access modular
forms database, and the PARI development team for the \texttt{lfun}
module which made all numerical claims of this paper accessible.
Computations were performed locally and on Vast.AI cloud
infrastructure. Hypothesis generation and preliminary literature
analysis were assisted by multi-LLM tooling (Claude Opus, DeepSeek);
all mathematical claims and all citation attributions appearing in this
paper were independently verified through PARI/GP computation, LMFDB
lookup, and ar\textsc{X}iv-API confirmation before inclusion. Several
earlier mis-attributions caught during this verification process
(notably an attempt to credit Conjecture~M183 to a fictitious
Kim--Lim 2017 paper, and a separate
attempt to draw a Darmon--Vonk connection that does not exist in the
literature) have been excised from the final text.

%% ============================================================
%% BIBLIOGRAPHY
%% ============================================================

\begin{thebibliography}{99}

%% ---- Verified arXiv references (post-1991, ID confirmed and
%% ---- title/author cross-checked via verify-arxiv.py + Sonnet body-scan) ----

\bibitem{BuyukbodukLei2016}
  K.~B\"uy\"ukboduk and A.~Lei,
  \emph{Anticyclotomic $p$-ordinary Iwasawa theory of elliptic modular forms},
  arXiv:1602.07508, 2016.

\bibitem{Kriz2018}
  D.~Kriz,
  \emph{A new $p$-adic Maass--Shimura operator and supersingular Rankin--Selberg $p$-adic $L$-functions},
  arXiv:1805.03605, 2018.

\bibitem{KannoWatari2020}
  K.~Kanno and T.~Watari,
  \emph{$W=0$ complex structure moduli stabilization on CM-type K3 $\times$ K3 orbifolds: arithmetic, geometry and particle physics},
  arXiv:2012.01111, 2020.

\bibitem{DarmonVonk2021}
  H.~Darmon and J.~Vonk,
  \emph{Singular moduli for real quadratic fields: a rigid analytic approach},
  arXiv:2511.08159 (companion: Duke Math. J., 2021).

\bibitem{ConnesMarcolliRamachandran2005}
  A.~Connes, M.~Marcolli and N.~Ramachandran,
  \emph{KMS states and complex multiplication},
  Selecta Math.\ \textbf{11} (2005), 325--347;
  arXiv:math/0501424.

\bibitem{Moore1998}
  G.~Moore,
  \emph{Arithmetic and attractors},
  arXiv:hep-th/9807087, 1998.

\bibitem{ChamseddineConnesMarcolli2007}
  A.~H. Chamseddine, A.~Connes and M.~Marcolli,
  \emph{Gravity and the standard model with neutrino mixing},
  Adv.\ Theor.\ Math.\ Phys.\ \textbf{11} (2007), 991--1089;
  arXiv:hep-th/0610241.

\bibitem{DabrowskiMukhopadhyayPozar2025}
  L.~D\k a}browski, B.~Mukhopadhyay and T.~Po{\v z}ar,
  \emph{Spectral torsion of the internal noncommutative geometry of the Standard Model},
  arXiv:2511.08159, 2025.

\bibitem{LMFDB2026}
  The LMFDB Collaboration,
  \emph{The $L$-functions and Modular Forms Database},
  \url{https://www.lmfdb.org}, accessed 2026.

\bibitem{PARI215}
  The PARI Group,
  \emph{PARI/GP, version 2.15.4},
  Univ.\ Bordeaux, 2024,
  \url{https://pari.math.u-bordeaux.fr/}.

%% ---- CITE-NEEDED entries (pre-arXiv classical literature) ----

\bibitem{Damerell1971}
  R.~M.~Damerell,
  \emph{$L$-functions of elliptic curves with complex multiplication, I},
  Acta Arith.\ \textbf{17} (1970), 287--301;
  \emph{II}, Acta Arith.\ \textbf{19} (1971), 311--317.
  \texttt{[CITE-NEEDED::pre-arXiv]}

\bibitem{Yager1982}
  R.~I.~Yager,
  \emph{On two variable $p$-adic $L$ functions},
  Ann.\ of Math.\ \textbf{115} (1982), 411--449.
  \texttt{[CITE-NEEDED::pre-arXiv; the explicit
   ``$p$-inert'' case of M183 is not covered by the published Yager theorem]}

\bibitem{GoldsteinSchappacher1981}
  C.~Goldstein and N.~Schappacher,
  \emph{S\'eries d'Eisenstein et fonctions $L$ de courbes elliptiques \`a multiplication complexe},
  J.\ Reine Angew.\ Math.\ \textbf{327} (1981), 184--218.
  \texttt{[CITE-NEEDED::pre-arXiv;
   contains Th\'eor\`eme~5.1 used for Cl$(K)$-orbit Corollary]}

\bibitem{ChowlaSelberg1949}
  S.~Chowla and A.~Selberg,
  \emph{On Epstein's zeta function (I)},
  Proc.\ Nat.\ Acad.\ Sci.\ USA \textbf{35} (1949), 371--374.
  \texttt{[CITE-NEEDED::pre-arXiv]}

\bibitem{CITE-NEEDED-PS1971}
  I.~I.~Pjatecki\u{\i}-\v{S}apiro,
  \emph{Interrelations between the Tate conjecture and the Hodge conjecture for three-dimensional varieties} (in Russian),
  Izv.\ Akad.\ Nauk SSSR Ser.\ Mat.\ \textbf{35} (1971), 1026--1038.
  \texttt{[CITE-NEEDED::pre-arXiv]}

\bibitem{CITE-NEEDED-Tankeev1990}
  S.~G.~Tankeev,
  \emph{Surfaces of type K3 over number fields and the Mumford--Tate conjecture},
  Izv.\ Akad.\ Nauk SSSR Ser.\ Mat.\ \textbf{54} (1990), no.~4, 846--861;
  English transl., Math.\ USSR-Izv.\ \textbf{37} (1991).
  \texttt{[CITE-NEEDED::pre-arXiv]}

\bibitem{CITE-NEEDED-Shimura1971}
  G.~Shimura,
  \emph{Introduction to the Arithmetic Theory of Automorphic Functions},
  Princeton Univ.\ Press, 1971.
  \texttt{[CITE-NEEDED::standard monograph]}

\bibitem{CITE-NEEDED-Schappacher1988}
  N.~Schappacher,
  \emph{Periods of Hecke Characters},
  Lecture Notes in Mathematics \textbf{1301}, Springer, 1988.
  \texttt{[CITE-NEEDED::standard monograph]}

\bibitem{CITE-NEEDED-Katz1978}
  N.~M.~Katz,
  \emph{$p$-adic $L$-functions for CM fields},
  Invent.\ Math.\ \textbf{49} (1978), 199--297.
  \texttt{[CITE-NEEDED::pre-arXiv]}

\bibitem{CITE-NEEDED-CoatesWiles1977}
  J.~Coates and A.~Wiles,
  \emph{On the conjecture of Birch and Swinnerton-Dyer},
  Invent.\ Math.\ \textbf{39} (1977), 223--251.
  \texttt{[CITE-NEEDED::pre-arXiv]}

\bibitem{CITE-NEEDED-AtkinLehner1970}
  A.~O.~L.~Atkin and J.~Lehner,
  \emph{Hecke operators on $\Gamma_0(m)$},
  Math.\ Ann.\ \textbf{185} (1970), 134--160.
  \texttt{[CITE-NEEDED::pre-arXiv]}

\bibitem{CITE-NEEDED-Robert1973}
  G.~Robert,
  \emph{Unit\'es elliptiques}, Bull.\ Soc.\ Math.\ France, M\'em.\ \textbf{36}, 1973.
  \texttt{[CITE-NEEDED::pre-arXiv; classical Eisenstein--Kronecker treatise]}

\bibitem{CITE-NEEDED-BannaiKobayashi2011}
  K.~Bannai and S.~Kobayashi,
  \emph{Algebraic theta functions and the $p$-adic interpolation of Eisenstein--Kronecker numbers},
  J.\ Reine Angew.\ Math.\ \textbf{660} (2011), 1--67.
  \texttt{[CITE-NEEDED::Crelle journal verification]}

\bibitem{CITE-NEEDED-Borel1981}
  A.~Borel,
  \emph{Commensurability classes and volumes of hyperbolic 3-manifolds},
  Ann.\ Scuola Norm.\ Sup.\ Pisa Cl.\ Sci.\ (4) \textbf{8} (1981), 1--33.
  \texttt{[CITE-NEEDED::pre-arXiv; Bianchi orbifold volumes]}

\bibitem{CITE-NEEDED-KennedyPendleton1985}
  A.~D.~Kennedy and B.~J.~Pendleton,
  \emph{Quasi-Monte Carlo simulation in lattice gauge theory},
  Phys.\ Lett.\ B \textbf{156} (1985), 393--399.
  \texttt{[CITE-NEEDED::pre-arXiv lattice physics reference]}

\end{thebibliography}

%% ============================================================
\appendix
%% ============================================================

\section{Reproducible PARI/GP script}\label{sec:pari}

The following PARI/GP~2.15.4 script computes $\alpha_2(D)$ for any
imaginary quadratic field $K = \Q(\sqrt{D})$ of class number~$1$
using the \texttt{lfun} module (analytic continuation through the
functional equation) and \texttt{bestappr} for rational recognition.
Precision is set to $150$ decimal digits; increase \texttt{prec} for
larger discriminants.

\begin{verbatim}
/* alpha2_cm.gp  --  PARI/GP 2.15.4                                */
/* Compute alpha_2 = L(f,2)*Pi^2/Omega^4 for CM weight-5 newforms. */
/* IMPORTANT: do NOT use partial sums Sum(a_n / n^2) at s=2;       */
/* the series converges absolutely only for Re(s) > 3, so partial */
/* sums oscillate inside the critical strip. Use lfun for analytic */
/* continuation.                                                    */

alpha2_cm(D, N, prec=150) = {
  my(H, ff, f, L, val, omega, alpha, rat);
  localprec(prec);

  \\ 1. Find weight-5 newform of level N with CM by Q(sqrt(D))
  H  = mfinit([N, 5, Mod(kronecker(D,x), x^2-1)], 1);
  ff = mfeigenbasis(H);
  f  = [];
  for (i = 1, #ff,
    if (mfisCM(ff[i]) && mfcmfield(ff[i]) == D,
      f = ff[i]; break));
  if (#f == 0, error("No CM form found for D=", D, " N=", N));

  \\ 2. L-value L(f,2) via analytic continuation
  L   = lfunmf(f);
  val = lfun(L, 2);

  \\ 3. Chowla-Selberg period Omega_D
  omega = chowla_selberg(D);

  \\ 4. alpha_2 = L(f,2)*Pi^2/Omega^4
  alpha = val * Pi^2 / omega^4;

  \\ 5. Rational recognition
  rat = bestappr(alpha, 10^(prec - 10));
  printf("D=%d  alpha_2 = %s  denom=%d\n", D, rat, denominator(rat));
  return(rat);
}

chowla_selberg(D) = {
  my(absD=abs(D), chi=kronecker(D,), w, prod=1.0);
  w = if(D==-3, 6, D==-4, 4, 2);
  forstep(a=1, absD-1, 1,
    if(gcd(a, absD)==1,
      prod *= gamma(a/absD)^(w/4 * chi(a))));
  return(prod / sqrt(2*Pi));
}

\\ Example invocations (matching Theorem 2.1):
\p 150
alpha2_cm(-4,  4)    \\ expects 1/12         (LMFDB 4.5.b.a)
alpha2_cm(-3,  12)   \\ expects 3/8          (LMFDB 12.5.b.a)
alpha2_cm(-7,  7)    \\ expects 28/3
alpha2_cm(-11, 11)   \\ expects 9/11  (first curve)
alpha2_cm(-19, 19)   \\ expects 13/57 (first curve)
alpha2_cm(-43, 43)   \\ expects 214/129
alpha2_cm(-67, 67)   \\ expects 1519/201
alpha2_cm(-163,163)  \\ expects 196216792/3 (first eigenform only)
\end{verbatim}

\begin{remark}
For $D \in \{-11, -19\}$, two distinct elliptic curves over $\Q$ with
CM by $\mathcal{O}_{K_D}$ exist; rerun the script after substituting the
alternative curve to recover both values of the twin pair. For
$D = -163$, two rational eigenforms exist at level $163$; only the
first is covered here.
\end{remark}

\end{document}
